% ============================================================
%  Résumé
% ============================================================
\chapter*{Résumé}
\addcontentsline{toc}{chapter}{Résumé}
\markboth{Résumé}{Résumé}

Ce rapport présente la conception et l'implémentation d'un système complet de
représentation des connaissances et de raisonnement (\textsc{rcr}), développé
dans le cadre du module \textsc{rcr1} du Master~1 \textsc{sii} à l'\textsc{usthb}.

Le projet met en œuvre cinq formalismes classiques de la représentation des
connaissances, appliqués à deux domaines concrets : l'\textbf{industrie musicale
algérienne} (artistes, labels, albums, certifications) et le
\textbf{système judiciaire} (accusés, magistrats, infractions, tribunaux).

\medskip
\noindent Les formalismes implémentés sont :

\begin{itemize}
  \item \textbf{Logique Modale} — modèles de Kripke, opérateurs de nécessité
        ($\nec$) et de possibilité ($\pos$), évaluation sémantique dans des
        mondes possibles.
  \item \textbf{Logique des Défauts} — théories de Reiter $\Delta = (W, D)$,
        calcul d'extensions, raisonnement non-monotone.
  \item \textbf{Logique Classique} — logique propositionnelle en forme normale
        conjonctive (CNF) résolue par SAT4J, et logique des prédicats du premier
        ordre avec chaînage avant.
  \item \textbf{Logique de Description} — ontologies OWL~2 construites via
        l'OWL~API~4.5, raisonnées par le moteur HermiT~1.4.3 ; TBox/ABox,
        subsomption, disjonction, classification des individus.
  \item \textbf{Réseaux Sémantiques} — graphes de concepts avec héritage et
        exceptions via TweetyProject.
\end{itemize}

\medskip
\noindent L'ensemble est accessible via une \textbf{interface graphique JavaFX~17}
interactive organisée en tableau de bord, avec un panneau dédié à chaque
formalisme, un affichage des résultats du raisonneur en arrière-plan, et une
démonstration explicite de la non-monotonie pour les logiques qui la supportent.

Le code source est disponible sur GitHub :
\url{https://github.com/D-Arslan/knowledge-representation-reasoning}

\bigskip
\noindent\textbf{Mots-clés :} représentation des connaissances, raisonnement,
logique modale, logique des défauts, logique de description, réseaux sémantiques,
OWL, HermiT, TweetyProject, non-monotonie.

\newpage
